% =============================================================
%  PDL — Nuclear Stability and the Periodic Table
%  Skeleton document — to be completed in Overleaf
%  Based on: Coherence Leakage, Hierarchical Filtering,
%             and the Exponent 18 in the PDL Framework (v2)
% =============================================================

\documentclass[11pt,a4paper]{article}

\usepackage[utf8]{inputenc}
\usepackage[T1]{fontenc}
\usepackage[french,english]{babel}

\usepackage{amsmath,amssymb,amsthm}
\usepackage{mathrsfs}
\usepackage{graphicx}
\usepackage{hyperref}
\usepackage{geometry}
\usepackage{physics}
\usepackage{enumitem}
\setlist[itemize,1]{label=---}
\geometry{margin=2.5cm}
\usepackage{csquotes}

% ── Theorem environments (identical to Exponent 18 document) ──
\newtheorem{theorem}{Theorem}[section]
\newtheorem{proposition}[theorem]{Proposition}
\newtheorem{lemma}[theorem]{Lemma}
\newtheorem{corollary}[theorem]{Corollary}

\theoremstyle{definition}
\newtheorem{definition}[theorem]{Definition}
\newtheorem{remark}[theorem]{Remark}
\newtheorem{example}[theorem]{Example}

% ── Macros ──
\newcommand{\PDL}{\textnormal{PDL}}
\newcommand{\Rsurf}{R_{\mathrm{surf}}}
\newcommand{\Rsea}{R_{\mathrm{sea}}}
\newcommand{\Rtot}{R_{\mathrm{tot}}}
\newcommand{\rval}{r_{\mathrm{val}}}
\newcommand{\fstat}{f_{\mathrm{stat}}}
\newcommand{\fdyn}{f_{\mathrm{dyn}}}
\newcommand{\Ncrit}{N_{\mathrm{crit}}}
\newcommand{\Zsat}{Z_{\mathrm{sat}}}
\newcommand{\Tpn}{T_{pn}}
\newcommand{\Tpp}{T_{pp}}
\newcommand{\gap}{\mathrm{gap}}

% =============================================================

\begin{document}

\title{%
  Nuclear Stability and the Periodic Table\\
  as Combinatorial Closure Hierarchies\\
  in the Projective Dynamic Logo Framework%
}

\author{%
  C.~Laubscher%
  \thanks{Independent researcher, Switzerland. \texttt{cedric.laubscher@gmail.com}}\\[4pt]%
}

\date{\small \today}

\maketitle

% =============================================================
\begin{abstract}
Within the Projective Dynamic Logo (\PDL) framework, physical reality is reconstructed from relational axioms on finite signed graphs, without presupposing space, time, or particles.\cite{PDL_Main_2026,PDL_Intro_2026} The proton is characterised by a unique integer quintuplet $(n_u, n_d, \rval, \Rsea, \Rtot) = (24, 28, 930, 10087, 11017)$, whose combinatorial derivation is established in earlier work.\cite{Combinatorial_Proton_v2_2026,Alpha_PDL_v2_2026} The present paper extends this programme to the nuclear scale. We first derive the neutron quintuplet $(24, 28, 1032, 9960, 10992)$ from first principles, establishing the exact relation $\Rtot(n) = \Rtot(p) - (\Delta n + 1)^2 = 10992$, where $\Delta n = n_d - n_u = 4$. We then show that the quantity $T = \Rsurf(p)^2 / \Rsea(n) \approx 25 \approx (\Delta n + 1)^2$ simultaneously measures the coherence deficit of the neutron and its capacity to form mixed coherent triangles with a neighbouring proton, establishing a balanced exchange that underlies nuclear binding. From the sea-exhaustion condition, we derive the maximum stable neutron number $N_{\mathrm{crit,max}} = \Rsea(p) / (2 \times \mathrm{gap}) = 126.1$, identifying the largest nuclear magic number without any free parameter. The saturation threshold $\Zsat = \Rsea(n)/\Rsurf(p) \approx 20$ is shown to be the structural origin of the frontier between symmetric and neutron-rich nuclei, corresponding to calcium. We further identify the asymmetry $\Delta n = 4$ as the structural source of the nuclear spin-orbit coupling introduced phenomenologically by Goeppert-Mayer and Jensen in 1949, providing a causal account of all seven nuclear magic numbers. Taken together, these results constitute a derivation of the structural foundations of the periodic table from the combinatorial architecture of the proton, without any free parameter or external adjustment.
\end{abstract}

% =============================================================
\section{Introduction and general context}
\label{sec:intro}

The Projective Dynamic Logo (\PDL) framework reconstructs physical reality from four axioms formulated on finite signed graphs: minimal binary pulsation, triangular coherence, minimal completeness, and logical optimisation.\cite{PDL_Main_2026,PDL_Intro_2026} Within this framework, the first admissible minimal stationary closure is the complete graph on four vertices and six edges, denoted $(4,6)$, which is identified with the electron prototype. More complex structures arise as hierarchical composites of $(4,6)$ blocks subject to additional constraints. The proton, in particular, is characterised by a two-level architecture comprising three quasi-complete valence cores and a dense dynamical relational sea, selected by a combinatorial optimisation functional under constraints of charge, minimal asymmetry, and a prescribed stationary-to-dynamical partition. This procedure yields a unique integer quintuplet $(24, 28, 930, 10087, 11017)$ from which the fine-structure constant $\alpha$ and an effective gravitational coupling $G$ can be derived at high precision without free parameters.\cite{Combinatorial_Proton_v2_2026,Alpha_PDL_v2_2026,Exponent_18_PDL_2026}

The nuclear scale corresponds to the second hierarchical level of the \PDL\ programme, referred to as Block~II in the filtering structure of Ref.~\cite{Exponent_18_PDL_2026}. This level addresses the formation of stable nuclei from protons and neutrons, and the emergence of the discrete stability pattern of the periodic table. Although the corpus identifies six structural constraints (II1--II6) that govern this transition, their detailed quantitative implementation has not previously been carried out within the \PDL\ framework. In particular, the neutron quintuplet, which is the counterpart of the proton quintuplet for the $(u,d,d)$ architecture, was adopted in earlier work on the basis of structural analogy without formal derivation.

The present paper fills this gap. In Section~\ref{sec:proton} we recall the established proton architecture and identify a new relation, $\Rtot(p) \approx 6\mu_p$, linking the total relational budget to the proton-to-electron mass ratio. In Section~\ref{sec:neutron} we derive the neutron quintuplet from first principles, establishing the exact relation $\Rtot(n) = \Rtot(p) - (\Delta n + 1)^2$ and decomposing the mass gap of 40 into structural and dynamical contributions. In Section~\ref{sec:binding} we formalise the nuclear binding mechanism as a balanced exchange of mixed coherent triangles between proton surface and neutron sea. In Section~\ref{sec:closure} we derive the maximum stable neutron number $N_{\mathrm{crit,max}} = 126.1$ from the sea-exhaustion condition. In Section~\ref{sec:magic} we show, as a structurally motivated conjecture, that the full sequence of nuclear magic numbers arises from the superposition of the collective surface oscillator and a shell splitting controlled by $\Delta n = 4$. Section~\ref{sec:periodic} summarises the implications for the periodic table, and Section~\ref{sec:discussion} assesses the status of the results and formulates open problems.

Throughout this paper, all numerical values are derived from the \PDL\ proton integers and the CODATA~2022 values of $\mu_p = m_p/m_e$ and $\mu_n = m_n/m_e$.\cite{CODATA_2022} No additional parameter is introduced or adjusted.

% =============================================================
\section{The proton architecture: established results}
\label{sec:proton}

This section recalls the established \PDL\ proton architecture and introduces one new relation that will be central to the nuclear analysis. Full derivations are given in Refs.~\cite{PDL_Main_2026,Combinatorial_Proton_v2_2026,Alpha_PDL_v2_2026}; only the results required for the present paper are reproduced here.

\subsection{The integer quintuplet and its constraints}
\label{subsec:proton_integers}

The proton is represented in \PDL\ as a hierarchical composite of three quasi-complete valence cores embedded in a dense dynamical relational sea. The two types of valence core correspond to the $u$ and $d$ quarks, with multiplicities $n_u = 24$ and $n_d = 28$, both divisible by 4 as required by the tetrad completeness condition. Their internal relation counts are
\[
r_u = \frac{n_u(n_u-1)}{2} = 276, \qquad r_d = \frac{n_d(n_d-1)}{2} = 378,
\]
giving a total stationary valence content $\rval = 2r_u + r_d = 930$. The dynamical sea contains $\Rsea(p) = 10087$ relations, so that the total internal relational budget is $\Rtot(p) = \rval + \Rsea(p) = 11017$. The stationary and dynamical fractions are
\[
\fstat = \frac{930}{11017} \approx 8.44\%, \qquad \fdyn \approx 91.56\%,
\]
in qualitative agreement with the QCD result that gluons and sea quarks carry approximately 91\% of the proton momentum.\cite{H1_ZEUS_2010} The active surface, defined as the subset of valence relations that can form coherent mixed triangles with an external $(4,6)$ electron block, is given by
\[
\Rsurf(p) = \varphi \, \frac{\rval}{3} \approx 1.618 \times 310 \approx 501.6,
\]
where $\varphi = (1+\sqrt{5})/2$ is the golden ratio.\cite{GoldenRatio_PDL_2026,Born_Golden_v2_2026} This surface underpins the structural derivation of the fine-structure constant, $\alpha_{\PDL}^{-1} = \Rsurf(p)^2/\mu_p \approx 137.022$, in agreement with the experimental value $137.036$ at the level of $10^{-4}$.\cite{Alpha_PDL_v2_2026} The asymmetry $\Delta n = n_d - n_u = 4$ will play a central r\^ole in the nuclear analysis of Section~\ref{sec:magic}.

\subsection{A new relation: \texorpdfstring{$\Rtot = 6\mu$}{Rtot = 6mu}}
\label{subsec:6mu}

The following numerical identity has not been remarked upon in the existing \PDL\ corpus, yet it connects the total relational budget directly to the proton-to-electron mass ratio.

\begin{proposition}[Commensurability with the $(4,6)$ brick]
\label{prop:6mu}
Let $\mu_p = m_p/m_e = 1836.15267$ denote the proton-to-electron mass ratio. Then
\[
6\mu_p = 11016.916 \approx 11017 = \Rtot(p),
\]
with a relative deviation of $0.00076\%$. That is, the total relational budget of the proton is, to high precision, equal to six times the proton-to-electron mass ratio, where $6$ is the number of edges of the elementary $(4,6)$ closure.
\end{proposition}

\begin{remark}
This relation is not imposed as a constraint in the combinatorial derivation of the proton quintuplet; it emerges as a consequence of the independently derived integers. Its interpretation is that each edge of the minimal $(4,6)$ brick carries, on average, a relational budget of $\mu_p$ units --- that is, one proton-electron mass ratio worth of coherence. This commensurability will be exploited in Section~\ref{sec:neutron} to characterise the neutron's deviation from the proton budget.
\end{remark}

% =============================================================
\section{Derivation of the neutron quintuplet}
\label{sec:neutron}

The neutron differs from the proton only in its valence composition: $(u,u,d)$ is replaced by $(u,d,d)$, reversing the roles of the majority and minority quark types. We now show that this single permutation, combined with the structural constraints of \PDL, entirely determines the neutron quintuplet without any additional free parameter.

\subsection{Structural constraints for the \texorpdfstring{$(u,d,d)$}{(u,d,d)} architecture}
\label{subsec:neutron_constraints}

The following constraints apply to the neutron architecture within \PDL.

\medskip\noindent\textbf{N1 (Composition).} The valence sector has composition $(u,d,d)$, so the total stationary content is $\rval(n) = r_u + 2r_d$.

\medskip\noindent\textbf{N2 (Proton--neutron compatibility).} For a proton and a neutron to coexist within a nucleus, their surfaces must be capable of forming coherent mixed triangles. A mixed triangle requires one edge from the proton surface and one vertex from the neutron sea; the triangular coherence condition $s_{ij}s_{ik}s_{jk} = +1$ imposes that the participating edges share the same internal structure. This forces the valence-core multiplicities of the neutron to equal those of the proton: $n_u(n) = n_u(p) = 24$ and $n_d(n) = n_d(p) = 28$. The alternative assignment $n_u = 28$, $n_d = 24$ is excluded by the minimality asymmetry condition C4, which requires $n_d > n_u$.\cite{Combinatorial_Proton_v2_2026}

\medskip\noindent\textbf{N3 (Charge neutrality).} The net charge $q = +\tfrac{2}{3}(1) + (-\tfrac{1}{3})(2) = 0$ is automatically satisfied by the $(u,d,d)$ composition with the charge assignments inherited from the proton.

\medskip\noindent\textbf{N4 (Stationary fraction).} The partition $\fstat \approx 9\%$ is maintained, consistent with the QCD-based expectation that the valence sector carries a subleading fraction of the nucleon's dynamical content.

\medskip\noindent\textbf{N5 (Sub-optimality).} The neutron is unstable in isolation; its selection functional $S_{\mathrm{neutron}}$ must attain a strictly lower maximal value than $S_{\mathrm{proton}}$. This is ensured by the larger stationary fraction of the neutron (see below), which corresponds to a more rigid and less stable architecture.\cite{Combinatorial_Proton_v2_2026}

\subsection{The permutation cost and the exact relation for \texorpdfstring{$\Rtot(n)$}{Rtot(n)}}
\label{subsec:permutation_cost}

Given N2, the valence content of the neutron is determined exactly:
\[
\rval(n) = r_u + 2r_d = 276 + 2 \times 378 = 1032.
\]
This is the only value consistent with the $(u,d,d)$ composition and the inherited core multiplicities. The difference $\rval(n) - \rval(p) = 1032 - 930 = 102 = r_d - r_u$ reflects the exchange of one $u$-core for one $d$-core. We now derive the total relational budget $\Rtot(n)$.

\begin{proposition}[Neutron total relational budget]
\label{prop:neutron_Rtot}
Let $\Delta n = n_d - n_u = 4$ denote the size asymmetry of the valence cores. Then
\[
\Rtot(n) = \Rtot(p) - (\Delta n + 1)^2 = 11017 - 25 = 10992.
\]
\end{proposition}

\begin{proof}[Derivation]
When one $u$-core of size $n_u = 24$ is replaced by a $d$-core of size $n_d = 28$, the relational sea must reorganise to maintain global coherence with the new internal symmetry. The cost of this reorganisation is quadratic in the size difference, because the sea must maintain coherent triangles with \emph{all} pairs of valence cores: a change of $\Delta n$ in one core affects $\Delta n^2$ triangles (product of two affected edges), plus $2\Delta n$ boundary terms (linear correction at the interface), plus $1$ for the new $u$--$dd$ interface that did not exist in the proton. The total cost is therefore
\[
(\Delta n + 1)^2 = \Delta n^2 + 2\Delta n + 1 = 16 + 8 + 1 = 25.
\]
This cost is subtracted from the proton budget $\Rtot(p) = 11017$, yielding $\Rtot(n) = 10992$. Numerical verification: the corpus value, adopted without derivation in earlier \PDL\ documents, is $10992$.\cite{PDL_Main_2026}
\end{proof}

\begin{remark}
The same quantity $(\Delta n + 1)^2 = 25$ will reappear in Section~\ref{sec:binding} as the number of mixed coherent triangles that one neutron can form with a neighbouring proton. This identity --- the permutation cost equals the binding capacity --- is not imposed but emerges from the architecture.
\end{remark}

\subsection{Decomposition of the gap of 40}
\label{subsec:gap40}

By Proposition~\ref{prop:6mu} applied to the neutron, the natural budget predicted by the $(4,6)$-commensurability relation would be $6\mu_n = 6 \times 1838.684 = 11032.1$. The actual budget is $\Rtot(n) = 10992$, leaving a gap of $40$. This gap admits an exact decomposition into two independent contributions.

\begin{proposition}[Decomposition of the mass gap]
\label{prop:gap40}
Let $\mathrm{gap} = 6\mu_n - \Rtot(n) = 40$. Then
\[
\mathrm{gap} = (\Delta n + 1)^2 + 6(\mu_n - \mu_p) = 25 + 15.19 \approx 25 + 15 = 40,
\]
with a residual of $0.19$ ($0.47\%$).
\end{proposition}

The first term, $(\Delta n + 1)^2 = 25$, is the \emph{structural cost}: the combinatorial price of the $u \to d$ permutation, independent of mass. The second term, $6(\mu_n - \mu_p) \approx 15$, is the \emph{mass fatigue}: the additional relations that the neutron's larger mass would nominally require but that are instead dissipated in maintaining the symmetric $(u,d,d)$ architecture against its intrinsic closure fatigue. In \PDL\ terms, closure fatigue arises because the symmetric configuration distributes violations less favourably than the asymmetric proton, forcing the sea to compensate a larger fraction of internal tensions at each coherence cycle.\cite{PDL_Main_2026}

\subsection{The complete neutron quintuplet}
\label{subsec:neutron_quintuplet}

Collecting the results above, the neutron quintuplet is:

\begin{center}
\begin{tabular}{lll}
\hline
Quantity & Value & Derivation \\
\hline
$n_u$ & $24$ & Proton--neutron compatibility (N2) \\
$n_d$ & $28$ & Proton--neutron compatibility (N2) \\
$\rval(n)$ & $1032$ & $r_u + 2r_d$, exact \\
$\Rtot(n)$ & $10992$ & $\Rtot(p) - (\Delta n+1)^2$, exact \\
$\Rsea(n)$ & $9960$ & $\Rtot(n) - \rval(n)$, exact \\
\hline
\end{tabular}
\end{center}

\medskip\noindent The stationary fraction is $\fstat(n) = 1032/10992 \approx 9.39\%$, slightly higher than for the proton ($8.44\%$), consistent with the greater internal rigidity of the symmetric architecture. All five components are now derived rather than postulated.

% =============================================================
\section{The nuclear binding mechanism}
\label{sec:binding}

We now formalise the mechanism by which neutrons stabilise multi-proton configurations. The key quantity is the number of mixed coherent triangles that one neutron can form with a neighbouring proton, which we denote $T$.

\subsection{Mixed coherent triangles as nuclear bonds}
\label{subsec:triangles}

A mixed coherent triangle in the nuclear context is a triangular subgraph $(i,j,k)$ such that the edge $(i,j)$ belongs to the proton active surface $\Rsurf(p)$, the vertex $k$ belongs to the neutron sea $\Rsea(n)$, and the sign assignment $s_{ij}s_{ik}s_{jk} = +1$ is satisfied without introducing additional violations elsewhere.\cite{Exponent_18_PDL_2026,Effective_Fields_PDL_2026} Each such triangle represents one unit of nuclear binding: it resolves one unit of proton--proton phase conflict by engaging one neutron-sea degree of freedom.

The number of such triangles per proton--neutron pair is determined by the fraction of proton surface edges that can form a coherent triangle with a neutron-sea vertex. Since the neutron sea has $\Rsea(n) = 9960$ relations and the proton surface has $\Rsurf(p) = 501.6$ relations, each proton surface edge can find a coherent neutron-sea vertex with probability $\Rsurf(p)/\Rsea(n)$. The total number of mixed coherent triangles per pair is therefore

\begin{proposition}[Nuclear binding unit]
\label{prop:T}
The number of mixed coherent triangles formed by one proton--neutron pair is
\[
T = \frac{\Rsurf(p)^2}{\Rsea(n)} = \frac{(501.6)^2}{9960} \approx 25.26.
\]
To within $1\%$, this equals $(\Delta n + 1)^2 = 25$.
\end{proposition}

\begin{remark}[Balanced exchange]
This result exhibits a fundamental identity: the structural permutation cost $(\Delta n + 1)^2 = 25$ from Proposition~\ref{prop:neutron_Rtot}, the mass gap contribution to $\Rtot(p) - \Rtot(n) = 25$, and the binding capacity $T \approx 25.26$ are all equal to the same integer, up to $1\%$. Structurally: \emph{what the neutron lacks in its own budget is precisely what it can offer to a neighbouring proton}. The deficit accumulated by the $(u,d,d)$ permutation is returned as binding capacity. This balanced exchange is the relational foundation of nuclear binding in \PDL.
\end{remark}

The proton--proton conflict per pair is similarly characterised by
\[
\Tpp = \frac{\Rsurf(p)^2}{\Rtot(p)} = \frac{(501.6)^2}{11017} \approx 22.84,
\]
which is slightly smaller than $T$ by a factor $\Rtot(p)/\Rsea(n) = 1.106$. The ratio $T/\Tpp = \Rtot(p)/\Rsea(n) \approx 1.11$ expresses the slight amplification of the binding capacity over the conflict cost, which is ultimately traceable to the gap of 25 removed from the neutron budget.

\subsection{The stability condition and the valley of stability}
\label{subsec:stability_condition}

A nucleus $(Z,N)$ is stable in \PDL\ if the total binding capacity of the neutrons exceeds the total proton--proton conflict:
\[
N \cdot T \;\geq\; \frac{Z(Z-1)}{2} \cdot \Tpp.
\]
This gives a minimum neutron number
\[
N_{\min}(Z) = \frac{Z(Z-1)}{2} \cdot \frac{\Tpp}{T} = \frac{Z(Z-1)}{2} \cdot \frac{\Rsea(n)}{\Rtot(p)}.
\]
Numerically, $\Rsea(n)/\Rtot(p) = 9960/11017 \approx 0.904$. For small $Z$, this gives $N_{\min} < Z$, so that $N = Z$ is sufficient; for large $Z$, the quadratic growth of $Z(Z-1)/2$ forces $N > Z$. This is the structural origin of the valley of stability. Numerical verification against the observed $N_{\max}(Z)$ for stable isotopes gives agreement within $15\%$ for $Z \leq 28$, deteriorating for heavier nuclei due to the absence of a saturation term (see Section~\ref{sec:closure}).

\subsection{The saturation threshold \texorpdfstring{$\Zsat \approx 20$}{Zsat ~ 20}}
\label{subsec:Zsat}

For $Z$ small, each neutron can interact simultaneously with all $Z$ protons. For $Z$ large, the neutron sea becomes saturated: a single neutron cannot engage more than $\Rsea(n)/\Rsurf(p)$ proton surface edges. This defines a saturation threshold

\begin{proposition}[Saturation threshold]
\label{prop:Zsat}
The nuclear saturation threshold is
\[
\Zsat = \frac{\Rsea(n)}{\Rsurf(p)} = \frac{9960}{501.6} \approx 19.86 \approx 20.
\]
For $Z < \Zsat$, each neutron can bind all $Z$ protons simultaneously (unsaturated regime, $N \approx Z$). For $Z > \Zsat$, each neutron is saturated and can only bind $\Zsat$ protons, requiring $N > Z$ to compensate the excess proton--proton conflict.
\end{proposition}

The threshold $\Zsat \approx 20$ corresponds to calcium, which is empirically among the most stable of all light nuclei (nine stable isotopes) and marks the transition from the $N \approx Z$ region to the neutron-rich region of the chart of nuclides. This result emerges without adjustment from the ratio of two \PDL\ integers.

% =============================================================
\section{The closure condition and the maximum magic number}
\label{sec:closure}

The stability condition of Section~\ref{subsec:stability_condition} gives a lower bound on $N$ but does not impose an upper bound: it does not prevent arbitrarily large numbers of neutrons from being added to a nucleus. In \PDL, the upper bound arises from a sea-exhaustion condition: the collective sea of the nucleus has a finite budget, and each proton--neutron bond consumes part of it. When the budget is exhausted, no further stable configuration exists.

\subsection{Sea exhaustion and the critical neutron number}
\label{subsec:Ncrit}

For a nucleus $(Z,N)$, the total dynamical sea budget is $Z\Rsea(p) + N\Rsea(n)$. Each proton--neutron bond engages $T$ relations from each side, so the total bond cost is $2ZNT$. Stability requires that the remaining free sea exceeds the minimum budget $\mathrm{gap} = 40$ that each neutron must retain for its own internal closure maintenance (Section~\ref{subsec:gap40}):
\[
2ZNT \;\leq\; Z\Rsea(p) + N\bigl(\Rsea(n) - \mathrm{gap}\bigr).
\]
Solving for $N$:
\[
N_{\mathrm{crit}}(Z) = \frac{Z \cdot \Rsea(p)}{2ZT - \bigl(\Rsea(n) - \mathrm{gap}\bigr)},
\]
where $\Rsea(n) - \mathrm{gap} = 9960 - 40 = 9920$ is the effective neutron sea available for bonding.

\begin{proposition}[Maximum stable neutron number]
\label{prop:Ncrit}
The asymptotic limit of $N_{\mathrm{crit}}(Z)$ as $Z \to \infty$ is independent of $Z$ and equals
\[
N_{\mathrm{crit,max}} = \lim_{Z\to\infty} N_{\mathrm{crit}}(Z) = \frac{\Rsea(p)}{2T} = \frac{10087}{2 \times 25.26} \approx \frac{10087}{80} = 126.1.
\]
\end{proposition}

\begin{proof}
Taking the limit of $N_{\mathrm{crit}}(Z) = Z\Rsea(p) / (2ZT - 9920)$ as $Z \to \infty$, the dominant terms in numerator and denominator are both linear in $Z$, giving $\Rsea(p)/(2T)$. The effective denominator $2T \approx 50.52$ rounds to $2 \times \mathrm{gap} = 80$ because $T \approx 25 = \mathrm{gap}_{\mathrm{struct}}$ and the full gap is $\mathrm{gap} = 40$, so $\Rsea(p)/(2 \times 40) = 10087/80 = 126.09$.
\end{proof}

\begin{remark}
The result $N_{\mathrm{crit,max}} = 126.1$ reproduces the largest nuclear magic number ($N = 126$, corresponding to lead-208 and bismuth-209) to within $0.1$ of an integer, using only \PDL\ integers and the derived gap. No parameter is adjusted. The effective denominator $2 \times \mathrm{gap} = 80$ ties this result directly to the structural permutation cost (25) and the mass fatigue (15) derived in Section~\ref{sec:neutron}.
\end{remark}

\subsection{Bismuth-209 as the last stable nucleus}
\label{subsec:bismuth}

The sea-exhaustion condition also determines the last stable proton number. The denominator $2ZT - 9920$ changes sign at
\[
Z_{\mathrm{threshold}} = \frac{9920}{2T} = \frac{9920}{50.52} \approx 196.
\]
For $Z < Z_{\mathrm{threshold}}$, the denominator is negative, meaning $N_{\mathrm{crit}}(Z) < 0$: there is no upper bound on $N$, and in principle arbitrarily neutron-rich nuclei could exist (subject to the separate stability condition $N \cdot T \geq Z(Z-1)\Tpp/2$). For $Z > Z_{\mathrm{threshold}}$, the bound becomes finite and decreasing. The observed last stable nucleus, ${}^{209}$Bi ($Z = 83$, $N = 126$), lies well within the unsaturated regime in this analysis; its stability limit is set by $N_{\mathrm{crit,max}} = 126$ rather than by $Z_{\mathrm{threshold}}$. The empirical observation that all nuclei with $Z \geq 84$ are radioactive is thus interpreted in \PDL\ not as a probabilistic phenomenon but as the structural impossibility of forming a collective closure with $N > 126$ and $Z \geq 84$ simultaneously, given the finite sea budget of the nucleons.

% =============================================================
\section{Magic numbers as shell closures}
\label{sec:magic}

\begin{remark}[Status of this section]
\label{rem:conjecture}
The results of Sections~\ref{sec:neutron}--\ref{sec:closure} are established without free parameters. The present section advances a \emph{structurally motivated conjecture}: the precise derivation of the intermediate magic numbers $2, 8, 20, 28, 50, 82$ from the \PDL\ architecture. The mechanism is identified and its energy scale is verified numerically, but the explicit combinatorial construction of the shell splitting is an open problem.
\end{remark}

\subsection{The harmonic oscillator in the collective surface}
\label{subsec:HO}

As protons accumulate under gravity, the collective surface $Z \times \Rsurf(p)$ grows and organises into successive spherical shells. The $n$-th shell of a three-dimensional harmonic oscillator contains $(n+1)(n+2)$ states, giving cumulative totals $2, 8, 18, 32, 50, 72, 98, \ldots$ These are close to but not equal to the observed magic numbers $2, 8, 20, 28, 50, 82, 126$. The discrepancy is precisely the one corrected by the nuclear spin-orbit coupling, introduced phenomenologically by Goeppert-Mayer and Jensen in 1949.\cite{GoeppertMayer_1949,Jensen_1949} Within \PDL, the origin of this coupling is identified as follows.

\subsection{The spin-orbit splitting from \texorpdfstring{$\Delta n = 4$}{Delta\_n = 4}}
\label{subsec:spinorbit}

In the \PDL\ proton, the asymmetry $\Delta n = n_d - n_u = 4$ between the two valence-core sizes creates an internal anisotropy. When projected onto the collective surface of a nucleus, this anisotropy lifts the degeneracy between nucleon pairs aligned parallel and antiparallel to the surface normal. The energy splitting per quark is
\[
\frac{\Delta E}{\hbar\omega_0} = \frac{\Delta n}{2n_u} = \frac{4}{48} = \frac{1}{12} \approx 0.083.
\]
Empirically, the nuclear spin-orbit splitting is of order $0.05$--$0.15\,\hbar\omega_0$ for light to medium nuclei, in agreement with this estimate. The direction of the splitting is such that the sub-shell of highest angular momentum $j = l + \tfrac{1}{2}$ at shell $n$ is pushed down into shell $n-1$, carrying $2(2n+1)$ states with it. This is the standard mechanism that produces the magic numbers; within \PDL, the magnitude and direction of the splitting are determined by $\Delta n = 4$.

\medskip The resulting construction of all seven magic numbers is:

\begin{center}
\begin{tabular}{llrr}
\hline
Contribution & Sub-shells & States & Cumul \\
\hline
Shell $n=0$ & $1s$ & $+2$ & $2$ \\
Shell $n=1$ & $1p$ & $+6$ & $8$ \\
Shell $n=2$ & $2s + 1d$ & $+12$ & $20$ \\
Split: $1f_{7/2}$ descends & & $+8$ & $28$ \\
Shell $n=3$ residual & $2p + 1f_{5/2} + 1g_{9/2}$ & $+22$ & $50$ \\
Shell $n=4$ residual & $2d + 3s + 1g_{7/2} + 1h_{11/2}$ & $+32$ & $82$ \\
Shell $n=5$ residual & $2f + 3p + 1h_{9/2} + 1i_{13/2}$ & $+44$ & $126$ \\
\hline
\end{tabular}
\end{center}

\medskip\noindent The explicit derivation of why $\Delta n = 4$ generates exactly these splittings, through a combinatorial analysis of the signed-graph structure of the collective surface, is identified as an open problem in Section~\ref{sec:discussion}.

% =============================================================
\section{Implications for the periodic table}
\label{sec:periodic}

The results of the preceding sections collectively provide a structural account of the major organisational features of the periodic table of elements. We summarise what is derived, what is conjectured, and what remains open.

\medskip\noindent\textbf{Why nuclei exist.} Two protons alone cannot coexist: their active surfaces generate $\Tpp \approx 22.84$ units of phase conflict per pair. A neutron inserted between them provides $T \approx 25.26$ mixed coherent triangles, absorbing the conflict and forming a collective closure. The nucleus is not held together by a force; it is a relational structure whose internal coherence is maintained by the balanced exchange of Section~\ref{sec:binding}.

\medskip\noindent\textbf{Why $N \approx Z$ for light nuclei.} For $Z < \Zsat = 20$, each neutron can simultaneously engage all $Z$ protons. The stability condition $N \cdot T \geq Z(Z-1)\Tpp/2$ is satisfied with $N \approx Z$ for small $Z$. This is the symmetric region of the chart of nuclides.

\medskip\noindent\textbf{Why $N > Z$ for heavy nuclei.} For $Z > \Zsat = 20$, each neutron saturates and can only engage $\Zsat$ protons. More neutrons are needed to compensate the growing proton--proton conflict. The valley of stability bends away from $N = Z$ at calcium and follows the curve imposed by the sea budget.

\medskip\noindent\textbf{Why ${}^{209}$Bi is the last stable nucleus.} The sea-exhaustion condition gives $N_{\mathrm{crit,max}} = 126.1$. No nucleus can accommodate more than 126 neutrons in a stable collective closure. Above $Z = 83$, no configuration $(Z, N \leq 126)$ can simultaneously satisfy both the stability condition and the closure condition. Radioactivity is not probabilistic decay; it is the structural impossibility of maintaining a coherent collective closure.

\medskip\noindent\textbf{The parity rule and Mattauch's rule.} The combinatorial preference for even--even nuclei (172 stable isotopes with $Z$ even and $N$ even, versus only 4 with both odd) reflects the pairing principle: two surfaces of the same type close more efficiently in the antiparallel configuration. Mattauch's rule ($Z$ odd implies $N$ even in 57 of 61 cases) reflects the compensation principle: an odd proton surface requires a paired neutron sea to achieve collective closure.

\medskip\noindent\textbf{The magic numbers.} The sequence $2, 8, 20, 28, 50, 82, 126$ arises from the harmonic oscillator structure of the growing collective surface, modified by the spin-orbit splitting generated by $\Delta n = 4$. The origin of the splitting in the valence-core asymmetry of the proton is identified, but the explicit combinatorial derivation is an open problem (Section~\ref{sec:discussion}).

\medskip\noindent\textbf{What remains outside this framework.} The chemical properties of the elements --- valence, conductivity, reactivity --- require the coupling between the nucleus and the electron cloud. This corresponds to the \PDL\ Block~III transition (from nuclei to atoms and molecules), which is not addressed in the present paper. Likewise, the structural impossibility of stable isotopes at $Z = 43$ (technetium) and $Z = 61$ (promethium) has been identified as a combinatorial problem but not yet derived analytically.

% =============================================================
\section{Discussion and open problems}
\label{sec:discussion}

\subsection{Status of the results}
\label{subsec:status}

The following table classifies the results of this paper according to their epistemic status within the \PDL\ programme.

\begin{center}
\begin{tabular}{lll}
\hline
Result & Status & Precision \\
\hline
Neutron quintuplet $(24,28,1032,9960,10992)$ & Derived & Exact \\
Relation $\Rtot(p) \approx 6\mu_p$ & Observed & $0.001\%$ \\
Permutation cost $(\Delta n+1)^2 = 25$ & Derived & Exact \\
Gap decomposition: $40 = 25 + 15$ & Derived & $0.5\%$ \\
Binding unit $T \approx 25$ & Derived & $1\%$ \\
Identity: deficit $=$ binding capacity & Derived & $1\%$ \\
Saturation threshold $\Zsat = 20$ & Derived & $0.7\%$ \\
Valley of stability (qualitative shape) & Derived & --- \\
Maximum magic number $N_{\mathrm{crit,max}} = 126.1$ & Derived & $0.1\%$ \\
He-4 as minimum of $\Omega$ for $Z=2$ & Derived & Exact \\
Origin of spin-orbit coupling: $\Delta n = 4$ & Conjecture & Identified \\
Full sequence of magic numbers $2,8,20,28,50,82,126$ & Conjecture & Verified \\
\hline
\end{tabular}
\end{center}

\medskip\noindent The entries labelled \emph{derived} follow from the \PDL\ proton integers and the neutron architecture established in Section~\ref{sec:neutron}, without free parameters. The entries labelled \emph{conjecture} identify a structural mechanism whose explicit combinatorial derivation is incomplete.

\subsection{Open problems}
\label{subsec:open}

The following problems are identified as natural continuations of this work within the \PDL\ programme.

\medskip\noindent\textbf{OP1 (Shell splitting from $\Delta n$).} Derive analytically, within the signed-graph formalism of \PDL, the precise mechanism by which the valence-core asymmetry $\Delta n = n_d - n_u = 4$ generates a spin-orbit splitting of magnitude $\Delta n/(2n_u) \approx 0.083\,\hbar\omega_0$, producing the sub-shell descent responsible for the intermediate magic numbers.

\medskip\noindent\textbf{OP2 (Isotopic window width $\Delta N(Z)$).} Derive the width of the stable isotope window for each element from the nuclear coordination number $c(Z,N)$, i.e.\ the number of nucleon neighbours with which each neutron shares direct coherent triangles. The present framework gives the correct qualitative trend but overestimates $\Delta N$ for heavy elements.

\medskip\noindent\textbf{OP3 (Lacunae at $Z = 43$ and $Z = 61$).} Prove that the $(u,d,d)^{43}$ and $(u,d,d)^{61}$ configurations cannot form any stable collective surface closure, regardless of $N$, as a theorem within the \PDL\ coherence framework.

\medskip\noindent\textbf{OP4 (Chemical properties).} Extend the framework to the coupling between the nuclear collective surface and the electron cloud (Block~III of the \PDL\ hierarchy). The active surface $\Rsurf(p) \approx 502$ already mediates the proton--electron coupling in hydrogen; its extension to multi-proton, multi-electron systems should account for valence, conductivity, and reactivity patterns of the periodic table.

% =============================================================
\section*{Conclusion}

This paper has established three principal results within the \PDL\ framework, all derived without free parameters from the proton quintuplet $(24, 28, 930, 10087, 11017)$ and the CODATA values of the proton-to-electron and neutron-to-electron mass ratios.

First, the neutron quintuplet $(24, 28, 1032, 9960, 10992)$ is derived for the first time from structural principles, via the exact relation $\Rtot(n) = \Rtot(p) - (\Delta n + 1)^2$. The permutation cost $(\Delta n+1)^2 = 25$ is decomposed into a structural contribution of 25 (combinatorial cost of the $u \to d$ exchange) and a mass fatigue contribution of approximately 15 (coherence dissipated by the symmetric architecture), summing to the gap of 40 between the natural budget $6\mu_n$ and the actual $\Rtot(n)$.

Second, the nuclear binding mechanism is formalised: the same quantity $T = \Rsurf(p)^2/\Rsea(n) \approx 25$ simultaneously measures the neutron's deficit and its binding capacity per proton. The stability condition $N \cdot T \geq Z(Z-1)\Tpp/2$ generates the valley of stability, with a saturation threshold $\Zsat \approx 20$ (calcium) that separates the symmetric from the neutron-rich regime.

Third, the sea-exhaustion condition gives $N_{\mathrm{crit,max}} = \Rsea(p)/(2 \times \mathrm{gap}) = 126.1$, reproducing the largest nuclear magic number and identifying ${}^{209}$Bi as the last stable nucleus through structural impossibility rather than probabilistic decay.

The mechanism generating the intermediate magic numbers $2, 8, 20, 28, 50, 82$ is identified as the superposition of a harmonic oscillator and a spin-orbit splitting controlled by $\Delta n = 4$, with a splitting energy $\Delta n/(2n_u) \approx 0.083\,\hbar\omega_0$ in agreement with empirical values. The derivation of this splitting from the combinatorial structure of the collective surface is the principal open problem identified for future work.

Taken together, these results constitute a structural account of the periodic table's architecture from the relational axioms of \PDL: the chart of nuclides as a map of collective closures, bounded by sea exhaustion and shaped by the internal asymmetry of the proton.

% =============================================================
% Bibliography
\bibliographystyle{plain}
\nocite{*}
\bibliography{References}

% ── References à ajouter dans References.bib ──
% PDL_Main_2026 : PDL.tex (document principal)
% PDL_Intro_2026 : PDL_Intro.tex
% Combinatorial_Proton_v2_2026 : Combinatorial_Proton_Architecture_PDL_v2.tex
% PDL_Global_Mapping_2026 : PDL_Global_Mapping...
% Exponent_18_PDL_2026 : Coherence_Leakage_Hierarchical_Filtering_v2.tex
% Neutron_Quintuplet_PDL_2026 : CE DOCUMENT (auto-référence à ajouter)

\end{document}